\documentclass{report}
\usepackage{amsmath,amssymb}
\setlength{\parindent}{0mm}
\setlength{\parskip}{1em}
\begin{document}
\begin{center}
\rule{6in}{1pt} \
{\large HyeongKae Park \\
{\bf Multigrid preconditioning for Monte Carlo solution of thermal radiative transfer problems}}

Los Alamos National Laboratory \\ P O Box 1663 T-3 MS-B216 \\ Los Alamos \\ NM 87545
\\
{\tt hkpark@lanl.gov}\\
Dana A. Knoll\\
Rick M. Rauenzahn\end{center}

Thermal radiative transfer problems appear in many different physical
applications such as combustion systems, inertial confinement fusion, and
astrophysics. We are developing a moment-based scale-bridging algorithm,
which utilizes the phase-space moments of the original transport equation
to accelerate the convergence of isotropic physics, such as
absorption-emission physics. In thermal radiative transfer problems, the
resulting low-order (LO), continuum system consists of coupled hyperbolic
equations and a ODE.

An efficient and accurate solver for the LO system is a key for
successful application of the moment-based, scale-bridging algorithm. The
LO system is solved using the Jacobian-free, Newton-Krylov (JFNK) method.
JFNK together with nonlinear elimination of the radiation flux and
material temperature equations, we can effectively reduce the LO system
to a scalar PDE for the radiation energy density. This scalar PDE takes a
form of an advection-diffusion equation that is amenable to effective
multigrid preconditioning.

However, when a Monte Carlo method is used to solve a (simplified)
high-order transport method, we encounter a difficulty due to the
presence of a stochastic noise in an advection (consistency) term. The
stochastic noise reduces the diagonal dominance of the preconditioning
system, which may lead to an eventual failure of multigrid
preconditioning. In this talk, we present effectiveness of multigrid
preconditioning in Monte Carlo thermal radiative transfer problems and
discuss a strength and weakness.


\end{document}
